\documentclass[11pt,a4paper]{article}
\usepackage{amsmath,amssymb,amsthm}
\usepackage[margin=2.5cm]{geometry}
\usepackage{hyperref}

\newtheorem{definition}{Definition}[section]
\newtheorem{theorem}[definition]{Theorem}
\newtheorem{proposition}[definition]{Proposition}
\newtheorem{remark}[definition]{Remark}
\newtheorem{conjecture}[definition]{Conjecture}

\title{Hyperseed Formalization:\\
  Tensor Logic for Bridging Neural and Symbolic AI}
\author{ProtomegaTron (automated formalization)}
\date{2026-09-18}

\begin{document}
\maketitle

\begin{abstract}
We formalize the intellectual structure of Ben Goertzel's Substack article
``Tensor Logic for Bridging Neural and Symbolic AI'' (2025-12-16) within
the Hyperseed ontology. The article describes how tensor logic bridges
symbolic reasoning and neural networks by representing logical databases
as sparse tensors and inference rules as tensor contractions. Key mappings:
logic-as-tensor as fiber-as-matrix, the integration problem as fiber-base
interface gap, tensor contraction as fiber transport, PLN truth values as
fiber density, unified substrate as fiber-base unification, and gradient
flow as differentiable fiber.
\end{abstract}

\tableofcontents

%% =================================================================
\section{Fiber-as-matrix: logic is linear algebra}
\label{sec:fiber}
%% =================================================================

\begin{theorem}[Logic as tensor operations --- claims 5--9, 26]
\label{thm:logic}
The fundamental observation: logical databases are sparse tensors, and
logical rules are tensor contractions. A relation $P(x,y)$ is a sparse
Boolean matrix $P_{xy}$. A rule like ``$P(x,y) \wedge P(y,z) \Rightarrow
G(x,z)$'' becomes matrix multiplication:
\[
  G_{xz} = H\left(\sum_y P_{xy} \cdot P_{yz}\right)
\]
where $H$ is a threshold function. This is exactly what GPUs excel at.

In Hyperseed terms, this is \emph{fiber-as-matrix}: the fiber (logical
structure --- relations, rules, inference chains) can be represented as
a tensor (matrix), enabling GPU-accelerated fiber operations:
\[
  F_{\text{logic}} \xrightarrow{\text{encode}} T_{\text{sparse}}
  \xrightarrow{\text{GPU}} T'_{\text{sparse}}
  \xrightarrow{\text{decode}} F'_{\text{logic}}
\]

The fiber's discrete relational structure maps to sparse tensor entries;
fiber transport (inference) maps to tensor contractions. This isn't a
notational trick --- it means logical reasoning can genuinely run on GPU
hardware at GPU speeds, using the same kernels that power deep learning.

The consequence: neural computations and logical inference can be seamlessly
mixed in the same computational graph. Both are tensor operations; the
distinction between ``neural'' and ``logical'' dissolves at the hardware
level.
\end{theorem}

%% =================================================================
\section{Fiber-base interface: the integration problem}
\label{sec:interface}
%% =================================================================

\begin{proposition}[Integration problem --- claims 1--4, 27]
\label{prop:integration}
Symbolic systems love discrete structures (graphs, trees, logical formulas).
Neural networks love dense matrices and tensors. When you build hybrid
systems, you get awkward translation layers and serialization bottlenecks.

In Hyperseed terms, this is a \emph{fiber-base interface gap}: neural
networks operate on the base (continuous tensor space) while logic operates
on the fiber (discrete relational structure). They represent the same
underlying cognitive structure in incompatible forms:
\[
  F_{\text{discrete}} \neq B_{\text{continuous}}
  \quad \text{but} \quad
  F_{\text{discrete}} \xrightarrow{\text{tensor logic}}
  T_{\text{compatible with } B}
\]

Tensor logic bridges the gap by representing the fiber in base-compatible
form (tensors). The fiber retains its logical structure but gains
base-compatible representation, eliminating the translation bottleneck.

The Hyperon project (MORK metagraph database, PeTTa/MM2 MeTTa interpreters,
PRIMUS cognitive architecture) needs this bridge: all components need
efficient neural-symbolic interoperation, which tensor logic provides.
\end{proposition}

%% =================================================================
\section{Fiber transport as tensor contraction}
\label{sec:transport}
%% =================================================================

\begin{theorem}[Tensor contraction = fiber transport --- claims 6--7, 28]
\label{thm:transport}
Inference rules as tensor contractions are fiber transport operations in
matrix form. Applying a rule ``if $P(x,y)$ then $Q(x)$'' is transporting
fiber content from the $P$-fiber to the $Q$-fiber via contraction:
\[
  Q_x = \sum_y P_{xy} \cdot w_{y} \qquad \text{(fiber transport)}
\]

The Einstein summation (einsum) notation provides a general language for
expressing arbitrary fiber transport patterns. Any inference rule ---
no matter how complex --- can be expressed as an einsum operation over
the relevant tensors.

This connects to the d-calculus: fiber transport in the d-calculus
(connection maps that move meaning along paths) corresponds to tensor
contractions that propagate truth values along inference paths. The
d-calculus provides the geometric interpretation; tensor logic provides
the computational implementation.
\end{theorem}

%% =================================================================
\section{Fiber density: PLN truth values}
\label{sec:pln}
%% =================================================================

\begin{proposition}[PLN as tensor logic --- claims 11, 15--18, 29]
\label{prop:pln}
PLN naturally fits tensor logic. The key extension: from Boolean (0/1)
to real-valued truth values for uncertain reasoning:
\[
  T_{xy} = (s_{xy}, c_{xy}) \in [0,1]^2
  \quad \text{(strength, confidence)}
\]

In Hyperseed terms, PLN's (strength, confidence) truth values correspond
to \emph{fiber density}: how much evidential weight the fiber carries at
each point. Boolean tensors give binary fiber (present/absent); real-valued
tensors give graded fiber (degrees of evidence).

PLN inference rules (deduction, induction, abduction, revision) can all
be expressed as tensor operations:
\begin{itemize}
\item \textbf{Deduction:} matrix multiplication with truth-value formulas.
\item \textbf{Revision:} weighted averaging of tensor entries.
\item \textbf{Induction/abduction:} transposed contractions with appropriate
  truth-value formulas.
\end{itemize}

This enables GPU-accelerated PLN reasoning --- orders of magnitude faster
than sequential symbolic PLN execution. For large knowledge bases with
millions of atoms, the speedup is transformative.
\end{proposition}

%% =================================================================
\section{Fiber-base unification: unified substrate}
\label{sec:unification}
%% =================================================================

\begin{theorem}[Unified substrate --- claims 19--22, 30--31]
\label{thm:unified}
Tensor logic provides a \emph{unified computational substrate}: both neural
and logical operations are tensor operations, running on the same hardware.

In Hyperseed terms, this is \emph{fiber-base unification}: the distinction
between fiber (logic) and base (neural) dissolves when both are represented
as tensor operations:
\[
  F_{\text{logic}} \cong T_{\text{logical}} \quad \text{and} \quad
  B_{\text{neural}} \cong T_{\text{neural}} \quad \implies \quad
  F \cup B \subset \text{Tensor}
\]

\textbf{Gradient flow = differentiable fiber:} If logical operations are
tensor operations, gradients can flow through logical inference. The fiber
becomes differentiable --- its structure can be optimized by gradient descent:
\[
  \nabla_\theta \mathcal{L}(\text{inference}(T_\theta))
  \quad \text{exists and is computable}
\]

This enables learning of logical rules from data: the fiber structure
itself becomes learnable, not just the fiber content.

\textbf{Attention as inference:} Transformer attention mechanisms can be
seen as logical inference in tensor form. The attention matrix $A_{ij}$
resembles a soft relational tensor; the attention operation resembles
tensor-logic inference. This connects transformers to logic at the
tensor level, suggesting that transformer ``reasoning'' is a form of
soft tensor-logic inference.

\textbf{Sparse optimization:} Logical tensors are typically very sparse
(most entities aren't related). Sparse tensor optimizations are crucial
for practical performance --- dense representations would waste memory
and compute on zero entries.
\end{theorem}

%% =================================================================
\section{Cross-references}
%% =================================================================

\begin{remark}[Connections to other formalizations]
\begin{itemize}
\item \textbf{Hyperseed v2} (2026-03-05): framework. Tensor logic
  provides a computational implementation of the Hyperseed fiber
  bundle --- the fiber's abstract structure gets a concrete tensor
  representation.
\item \textbf{Evidence Is to Logic What Energy Is to Physics} (2026-03-10):
  evidence conservation. Evidence conservation in quantales maps to
  tensor-norm conservation in tensor logic --- the total ``weight''
  of evidence is preserved under tensor contractions.
\item \textbf{OpenClaw: Amazing Hands for a Brain} (2026-02-03):
  brain improvement. Tensor logic is a fiber (brain) improvement,
  not a base (hands) improvement --- it makes the reasoning engine
  faster and more capable, not just the tool interface.
\item \textbf{Seeding RSI} (2026-07-28): architecture. Tensor logic
  as a key architectural component of Hyperon's reasoning engine.
\end{itemize}
\end{remark}

\begin{conjecture}[Tensor logic as universal fiber representation]
Conjecture: any fiber structure that can be expressed in first-order
(or higher-order) logic can be faithfully represented as a sparse
tensor, and any fiber transport operation can be expressed as a
tensor contraction.

If this conjecture holds, tensor logic provides a \emph{universal}
computational representation for Hyperseed fibers --- every fiber
structure has a tensor representation, and every fiber operation has
a tensor implementation. The Hyperseed ontology's abstract fiber
bundles become concrete tensor bundles, computable on GPU hardware.
\end{conjecture}

\end{document}
